\chapter{裸指针、节点和内存：Java 中该怎么想}

\section{指针不等于引用，但节点算法仍然一样}

你的 \code{clone-graph.cxx} 用 \code{Node*} 和 \code{unordered\_map<int, Node*>}。Java 中节点变量本质上是对象引用，图克隆的算法没有变化：

\begin{javacode}
class Node {
    int val;
    List<Node> neighbors = new ArrayList<>();
    Node(int val) { this.val = val; }
}

Node cloneGraph(Node start) {
    if (start == null) return null;

    Map<Node, Node> copies = new HashMap<>();
    Deque<Node> queue = new ArrayDeque<>();
    copies.put(start, new Node(start.val));
    queue.offer(start);

    while (!queue.isEmpty()) {
        Node current = queue.poll();
        for (Node next : current.neighbors) {
            if (!copies.containsKey(next)) {
                copies.put(next, new Node(next.val));
                queue.offer(next);
            }
            copies.get(current).neighbors.add(copies.get(next));
        }
    }
    return copies.get(start);
}
\end{javacode}

\begin{transfer}[用节点对象当 key]
你的原实现用 \code{node->val} 当哈希 key，隐含 ``节点值唯一'' 的前提。Java 直接用 \code{Map<Node, Node>}，把 ``节点身份'' 表达得更准确；这也是面试中更稳健的写法。
\end{transfer}

\section{LRU：先会手写，再知道标准库}

你在 \code{lru-cache.cxx} 维护 dummy head/tail、\code{prev/next} 和 \code{unordered\_map<int, Node*>}。这对于理解不变量非常好：哈希表定位节点，双向链表维护最近使用顺序。

生产 Java 的默认选择是 \code{LinkedHashMap}：

\begin{javacode}
class LRUCache extends LinkedHashMap<Integer, Integer> {
    private final int capacity;

    LRUCache(int capacity) {
        super(capacity, 0.75f, true); // true: access-order
        this.capacity = capacity;
    }

    public int get(int key) {
        return getOrDefault(key, -1);
    }

    public void put(int key, int value) {
        super.put(key, value);
    }

    @Override
    protected boolean removeEldestEntry(Map.Entry<Integer, Integer> eldest) {
        return size() > capacity;
    }
}
\end{javacode}

但 LeetCode 的 \code{LRUCache} 常要求你展示 \bigO{1} 设计。此时手写节点也完全合理，只需把 C++ 的 \code{Node*} 改成 \code{Node}，删除 \code{delete}，并让 \code{HashMap<Integer, Node>} 持有节点。

\section{GC 不是 ``不需要管理内存''}

Java 会回收不可达对象，但不会自动修复逻辑引用。以下情况仍会造成长期内存占用：
\begin{itemize}
  \item 全局 \code{static List} 持有每次测试数据；
  \item LRU 淘汰时只删链表、不删 \code{HashMap}；
  \item 图遍历的 \code{visited} 被复用而没有清空；
  \item Lambda/内部类无意捕获了大对象。
\end{itemize}

\begin{pitfall}[不要在 Java 里模拟 \code{delete}]
把节点从链表、Map、队列等所有可达路径断开就够了。调用 \code{System.gc()} 既不能保证回收，也不是算法题答案。
\end{pitfall}

\section{你该保留的 C++ 不变量意识}

FHQ Treap 的 \code{split/merge/pushup} 展示了优秀的 C++ 习惯：每个变换后恢复子树大小不变量。Java 写复杂树时仍应保留这种纪律：

\begin{javacode}
private void pull(Node node) {
    if (node != null) {
        node.size = size(node.left) + size(node.right) + node.count;
    }
}
\end{javacode}

语言换了，不变量、边界条件和析构前断链的意识不该换。
